\documentclass[11pt, a4paper]{article}

% bfw-style.tex — gemeinsame Style-Definitionen für alle Handouts/Aufgabenhefte
% Voraussetzung: Kompilation mit xelatex (wegen fontspec)

\usepackage[ngerman]{babel}
\usepackage{geometry}
\geometry{a4paper, top=2.2cm, bottom=2.2cm, left=2.3cm, right=2.3cm}

\usepackage{fontspec}
% Fallback-Kette: Source Sans Pro -> Calibri -> Segoe UI -> Latin Modern Sans
% (Latin Modern ist immer mit MiKTeX vorhanden)
\IfFontExistsTF{Source Sans Pro}{%
  \setmainfont{Source Sans Pro}[Ligatures=TeX]%
}{\IfFontExistsTF{Calibri}{%
  \setmainfont{Calibri}[Ligatures=TeX]%
}{\IfFontExistsTF{Segoe UI}{%
  \setmainfont{Segoe UI}[Ligatures=TeX]%
}{%
  \setmainfont{Latin Modern Sans}[Ligatures=TeX]%
}}}
\IfFontExistsTF{Source Code Pro}{%
  \setmonofont{Source Code Pro}[Scale=0.92]%
}{\IfFontExistsTF{Consolas}{%
  \setmonofont{Consolas}[Scale=0.92]%
}{%
  \setmonofont{Latin Modern Mono}[Scale=0.92]%
}}

\usepackage{xcolor}
% Farbpalette: hell, freundlich, modern
\definecolor{bfwPrimary}{HTML}{0EA5A4}    % Petrol/Türkis - Primär
\definecolor{bfwPrimaryDark}{HTML}{0F766E} % dunkler Petrol für Headings
\definecolor{bfwAccent}{HTML}{F59E0B}     % Amber - Akzent
\definecolor{bfwSuccess}{HTML}{10B981}    % Grün - Erfolg / Lösung
\definecolor{bfwWarn}{HTML}{F97316}       % Orange - Achtung
\definecolor{bfwInk}{HTML}{0F172A}        % fast Schwarz
\definecolor{bfwInkSoft}{HTML}{475569}    % gedämpftes Grau
\definecolor{bfwBgSoft}{HTML}{F8FAFC}     % off-white
\definecolor{bfwBgTeal}{HTML}{ECFEFF}     % helles Türkis
\definecolor{bfwBgAmber}{HTML}{FEF3C7}    % helles Gelb
\definecolor{bfwBgRose}{HTML}{FFE4E6}     % helles Rosé für Eselsbrücken

\usepackage[colorlinks=true, linkcolor=bfwPrimaryDark, urlcolor=bfwPrimaryDark]{hyperref}
\usepackage{enumitem}
\setlist[itemize]{leftmargin=*, itemsep=2pt, topsep=2pt}
\setlist[enumerate]{leftmargin=*, itemsep=2pt, topsep=2pt}

\usepackage[most]{tcolorbox}
\usepackage{booktabs}
\usepackage{tikz}
\usetikzlibrary{decorations.pathmorphing, positioning, shapes.geometric, arrows.meta}

% Merkkästchen — wichtige Definition / Kernpunkt
\newtcolorbox{merksatz}[1][]{%
  enhanced, breakable,
  colback=bfwBgTeal,
  colframe=bfwPrimary,
  boxrule=0.6pt,
  arc=4pt,
  left=10pt, right=10pt, top=8pt, bottom=8pt,
  fonttitle=\bfseries\color{bfwPrimaryDark},
  title={\faLightbulb\ Merksatz},
  #1
}

% Eselsbrücke — Mnemoniks
\newtcolorbox{eselsbruecke}[1][]{%
  enhanced, breakable,
  colback=bfwBgRose,
  colframe=bfwAccent,
  boxrule=0.6pt,
  arc=4pt,
  left=10pt, right=10pt, top=8pt, bottom=8pt,
  fonttitle=\bfseries\color{bfwWarn},
  title={Eselsbrücke},
  #1
}

% Beispiel-Box
\newtcolorbox{beispiel}[1][]{%
  enhanced, breakable,
  colback=bfwBgSoft,
  colframe=bfwInkSoft,
  boxrule=0.4pt,
  arc=3pt,
  left=10pt, right=10pt, top=8pt, bottom=8pt,
  fonttitle=\bfseries\color{bfwInk},
  title={Beispiel},
  #1
}

% Lösungs-Box (für Aufgabenhefte)
\newtcolorbox{loesung}[1][]{%
  enhanced, breakable,
  colback=bfwBgAmber!40,
  colframe=bfwSuccess,
  boxrule=0.6pt,
  arc=3pt,
  left=10pt, right=10pt, top=8pt, bottom=8pt,
  fonttitle=\bfseries\color{bfwSuccess!70!black},
  title={Lösung},
  #1
}

% Glossar-Eintrag
\newcommand{\glossareintrag}[2]{%
  \par\noindent\textbf{\color{bfwPrimaryDark}#1} \enspace #2\par\smallskip
}

% Section-Stil — etwas farbiger als Standard
\usepackage{titlesec}
\titleformat{\section}
  {\Large\bfseries\color{bfwPrimaryDark}}
  {\thesection}{0.6em}{}
\titleformat{\subsection}
  {\large\bfseries\color{bfwPrimary}}
  {\thesubsection}{0.5em}{}
\titleformat{\subsubsection}
  {\normalsize\bfseries\color{bfwInk}}
  {\thesubsubsection}{0.4em}{}

% TikZ-Stil "skizzenhaft" — etwas weicher als Rough.js, aber stabil
\tikzset{
  roughbox/.style = {
    draw=bfwPrimaryDark, thick, fill=bfwBgTeal,
    rounded corners=4pt, inner sep=6pt
  },
  roughaccent/.style = {
    draw=bfwAccent, thick, fill=bfwBgAmber,
    rounded corners=4pt, inner sep=6pt
  },
  roughline/.style = {
    draw=bfwInkSoft, thick, -{Stealth[length=2.5mm]}
  }
}

% Bullet-Symbole (bewusst kein FontAwesome — vermeidet Font-Installations-Probleme)
\newcommand{\faLightbulb}{$\bullet$}
\newcommand{\faBookmark}{$\bullet$}

% Header/Footer
\usepackage{fancyhdr}
\pagestyle{fancy}
\fancyhf{}
\renewcommand{\headrulewidth}{0pt}
\fancyfoot[L]{\small\color{bfwInkSoft}\thepage}
\fancyfoot[R]{\small\color{bfwInkSoft} BFW M-064 Beschaffung im Büromanagement}


\title{\vspace*{-1.5cm}\Huge\bfseries\color{bfwPrimaryDark}Aufgabenheft Tag~2\\[0.4em]
  \large\color{bfwInkSoft}\textnormal{Büroausstattung \& ergonomischer Bildschirmarbeitsplatz --- Übungen im IHK-Klausurformat}}
\author{\textsf{Büroprozesse gestalten und Arbeitsvorgänge organisieren \textperiodcentered{} M-067}\\
\textsf{\small\copyright\ Can Siebert\,\textperiodcentered\,2026}}
\date{Selbstlernmaterial \textperiodcentered{} 30.06.2026}

\begin{document}
\maketitle
\thispagestyle{empty}

\begin{merksatz}[title={\bfseries So nutzen Sie dieses Aufgabenheft}]
Bearbeiten Sie alle Aufgaben zunächst \emph{ohne} in die Lösungen zu schauen. Schreiben
Sie Ihre Antworten in vollständigen Sätzen --- die IHK-Prüfung verlangt formulierte
Begründungen, nicht bloße Stichworte. Beachten Sie die \emph{Operatoren} (nennen,
erläutern, beurteilen, begründen) --- sie geben den geforderten Antwort-Tiefegrad vor.
Erst danach öffnen Sie den Lösungsteil ab Seite~\pageref{loesungen} und vergleichen.
\end{merksatz}

\tableofcontents
\vspace{1em}
\hrule

\section{Aufgaben}

\subsection*{Aufgabe~1 --- Ergonomie verstehen (10 Punkte)}

\noindent\textbf{\color{bfwAccent}YouTube-Vorbereitung:}\enspace \href{https://www.youtube.com/results?search_query=Ergonomie+am+Arbeitsplatz+einfach+erklaert}{Ergonomie am Arbeitsplatz} (\url{https://www.youtube.com/results?search_query=Ergonomie+am+Arbeitsplatz+einfach+erklaert})

\medskip
Ihr Ausbildungsbetrieb plant eine Schulung zum Thema Gesundheit am Arbeitsplatz.

\begin{enumerate}[label=(\alph*)]
  \item \textbf{Erläutern} Sie, was man unter Ergonomie versteht. \hfill (3~P.)
  \item \textbf{Nennen} Sie das Ziel der Ergonomie für die arbeitenden Menschen. \hfill (3~P.)
  \item \textbf{Nennen} Sie drei Arbeitsmittel, die an die Körpermaße angepasst werden. \hfill (2~P.)
  \item \textbf{Beurteilen} Sie die Aussage: \emph{„Der Mensch muss sich an seinen Arbeitsplatz gewöhnen."} \hfill (2~P.)
\end{enumerate}

\subsection*{Aufgabe~2 --- Anforderungen an den Bürostuhl (12 Punkte)}

\noindent\textbf{\color{bfwAccent}YouTube-Vorbereitung:}\enspace \href{https://www.youtube.com/results?search_query=ergonomischer+Buerostuhl+richtig+einstellen}{Bürostuhl richtig einstellen} (\url{https://www.youtube.com/results?search_query=ergonomischer+Buerostuhl+richtig+einstellen})

\medskip
Die \emph{Lindner \& Partner GmbH} beschafft neue Bürodrehstühle.

\begin{enumerate}[label=(\alph*)]
  \item \textbf{Nennen} Sie die Norm, die Maße und Prüfkriterien für Bürodrehstühle regelt. \hfill (2~P.)
  \item \textbf{Nennen} Sie sechs Anforderungen an einen Bürodrehstuhl. \hfill (6~P.)
  \item \textbf{Erläutern} Sie, woran man erkennt, dass ein Stuhl korrekt eingestellt ist. \hfill (2~P.)
  \item \textbf{Begründen} Sie, warum mindestens 5 Rollen vorgeschrieben sind. \hfill (2~P.)
\end{enumerate}

\subsection*{Aufgabe~3 --- So sitzen Sie richtig (12 Punkte)}

\noindent\textbf{\color{bfwAccent}YouTube-Vorbereitung:}\enspace \href{https://www.youtube.com/results?search_query=Bildschirmarbeitsplatz+ergonomisch+einrichten}{Bildschirmarbeitsplatz einrichten} (\url{https://www.youtube.com/results?search_query=Bildschirmarbeitsplatz+ergonomisch+einrichten})

\medskip
\begin{enumerate}[label=(\alph*)]
  \item \textbf{Beschreiben} Sie, wo die oberste Bildschirmzeile liegen soll. \hfill (3~P.)
  \item \textbf{Nennen} Sie den empfohlenen Sichtabstand und die Position des Bildschirms zum Fenster. \hfill (3~P.)
  \item \textbf{Erläutern} Sie die korrekte Position von Tastatur und Maus. \hfill (3~P.)
  \item \textbf{Nennen} Sie die beiden Körperwinkel, die etwa 90\textdegree{} betragen sollen. \hfill (3~P.)
\end{enumerate}

\subsection*{Aufgabe~4 --- Mängel am Bildschirmarbeitsplatz (15 Punkte)}

\noindent\textbf{\color{bfwAccent}YouTube-Vorbereitung:}\enspace \href{https://www.youtube.com/results?search_query=Arbeitsplatz+ergonomisch+Fehler+vermeiden}{Ergonomie-Fehler vermeiden} (\url{https://www.youtube.com/results?search_query=Arbeitsplatz+ergonomisch+Fehler+vermeiden})

\medskip
Herr Below arbeitet täglich acht Stunden. Sein Laptop steht direkt auf der Tischplatte
(er schaut ständig nach unten), der Bildschirm steht frontal vor dem Fenster und spiegelt,
sein Stuhl ist nicht höhenverstellbar (die Füße baumeln), Tastatur und Maus nutzt er am
Laptop. Er sitzt den ganzen Tag in derselben Position und klagt über Nacken- und Rücken\-schmerzen.

\begin{enumerate}[label=(\alph*)]
  \item \textbf{Nennen} Sie fünf ergonomische Mängel an diesem Arbeitsplatz. \hfill (5~P.)
  \item \textbf{Schlagen} Sie zu jedem Mangel eine Verbesserungsmaßnahme \textbf{vor}. \hfill (5~P.)
  \item \textbf{Erläutern} Sie, warum ein Laptop im Dauereinsatz eine Dockingstation braucht. \hfill (3~P.)
  \item \textbf{Beurteilen} Sie, warum „den ganzen Tag gleich sitzen" gesundheitsschädlich ist. \hfill (2~P.)
\end{enumerate}

\subsection*{Aufgabe~5 --- Dynamisches Sitzen \& Sitzmechaniken (12 Punkte)}

\noindent\textbf{\color{bfwAccent}YouTube-Vorbereitung:}\enspace \href{https://www.youtube.com/results?search_query=dynamisches+Sitzen+Steh+Sitz+Dynamik}{Dynamisches Sitzen \& Steh-Sitz-Dynamik} (\url{https://www.youtube.com/results?search_query=dynamisches+Sitzen+Steh+Sitz+Dynamik})

\medskip
\begin{enumerate}[label=(\alph*)]
  \item \textbf{Erläutern} Sie das dynamische Sitzen und die drei Sitzpositionen. \hfill (3~P.)
  \item \textbf{Nennen} Sie die Faustregel der Schweizer Arbeitsmediziner (Sitzen/Stehen/Gehen). \hfill (3~P.)
  \item \textbf{Nennen} Sie, wie oft ein Haltungswechsel erfolgen und wie lang eine Stehphase höchstens sein soll. \hfill (2~P.)
  \item \textbf{Nennen} Sie die vier Sitzmechaniken und \textbf{erläutern} Sie eine davon. \hfill (4~P.)
\end{enumerate}

\subsection*{Aufgabe~6 --- Prüfsiegel, Tisch \& ArbStättV (15 Punkte)}

\noindent\textbf{\color{bfwAccent}YouTube-Vorbereitung:}\enspace \href{https://www.youtube.com/results?search_query=GS+Zeichen+CE+Kennzeichnung+Unterschied}{GS- \& CE-Zeichen} (\url{https://www.youtube.com/results?search_query=GS+Zeichen+CE+Kennzeichnung+Unterschied})

\medskip
\begin{enumerate}[label=(\alph*)]
  \item \textbf{Erläutern} Sie den Unterschied zwischen GS- und CE-Zeichen. \hfill (4~P.)
  \item \textbf{Nennen} Sie drei Anforderungen an den Arbeitstisch und die Höhe eines Steh-Sitz-Tischs. \hfill (4~P.)
  \item \textbf{Erläutern} Sie, wofür die Arbeitsstättenverordnung gilt und wo die Bildschirm\-anforderungen stehen. \hfill (3~P.)
  \item \textbf{Nennen} Sie vier Anforderungen der ArbStättV an Bildschirmarbeitsplätze. \hfill (4~P.)
\end{enumerate}

\vspace{1em}
\noindent\textbf{Punkte gesamt: 76}

\newpage

\section{Lösungen}\label{loesungen}

\subsection*{Lösung Aufgabe~1}
\begin{loesung}
\textbf{(a)} Ergonomie ist ein Teilgebiet der Arbeitswissenschaft. Sie befasst sich mit der
Abstimmung zwischen Mensch, Maschine und Arbeitswelt und dient dem präventiven
Gesundheitsschutz. Das Grundprinzip lautet: die Arbeit an den Menschen anpassen.

\textbf{(b)} Ziel ist es, die arbeitenden Menschen vor Berufskrankheiten, Arbeitsunfällen
und Monotonie zu bewahren.

\textbf{(c)} Beispiele: Stühle, Tische, Türgriffe (allgemein Arbeitsmittel, die an die Maße
des menschlichen Körpers angepasst werden).

\textbf{(d)} Die Aussage widerspricht der Ergonomie. Nicht der Mensch passt sich dem
Arbeitsplatz an, sondern der Arbeitsplatz wird an den Menschen angepasst --- nur so lassen
sich Belastungen, Beschwerden und Fehlzeiten vermeiden.
\end{loesung}

\subsection*{Lösung Aufgabe~2}
\begin{loesung}
\textbf{(a)} Die europäische Norm \textbf{DIN EN 1335}.

\textbf{(b)} Sechs Anforderungen (Auswahl): leichte Federung beim Setzen; keine scharfen
Kanten; gepolsterte, atmungsaktive Sitz- und Rückenlehnen; stufenlos verstellbare Sitzhöhe
42--50 cm; Sitztiefe mind.\ 38--44 cm; Sitzbreite mind.\ 40--48 cm; verstellbare
Rückenlehne (Breite mind.\ 36--48 cm); stand- und kippsicher; mind.\ 5 wegrollsichere Rollen;
Rollwiderstand passend zum Bodenbelag; höhenverstellbare Armlehnen.

\textbf{(c)} Bei aufrechter Haltung bilden Ober-/Unterarm und Ober-/Unterschenkel je einen
90\textdegree-Winkel, die Füße haben festen Bodenkontakt. Danach wird die Tischhöhe eingestellt.

\textbf{(d)} Mindestens 5 Rollen sorgen für Stand- und Kippsicherheit: Bei weniger Rollen
könnte der Stuhl beim Zurücklehnen oder Drehen umkippen --- ein Unfallrisiko.
\end{loesung}

\subsection*{Lösung Aufgabe~3}
\begin{loesung}
\textbf{(a)} Die oberste Bildschirmzeile soll \emph{leicht unterhalb der waagerechten
Sehachse} liegen, sodass der Blick leicht gesenkt ist.

\textbf{(b)} Sichtabstand mindestens 50 cm (praktisch 50--70 cm); der Bildschirm soll im
\emph{rechten Winkel zum Fenster} stehen, um Reflexionen und Blendungen zu vermeiden.

\textbf{(c)} Tastatur und Maus liegen in einer Ebene mit Ellenbogen und Handflächen; vor der
Tastatur muss Platz zum Auflegen der Handballen sein.

\textbf{(d)} 90\textdegree-Winkel zwischen Ober- und Unterarm sowie zwischen Ober- und
Unterschenkel; zusätzlich feste Fußauflage (ggf.\ Fußhocker).
\end{loesung}

\subsection*{Lösung Aufgabe~4}
\begin{loesung}
\textbf{(a)} Mängel (Auswahl): (1) Bildschirm zu niedrig, Blick zu stark nach unten;
(2) Bildschirm frontal vor Fenster $\to$ Spiegelungen/Blendung; (3) Stuhl nicht
höhenverstellbar, kein 90\textdegree-Winkel; (4) Füße ohne feste Auflage; (5) Tastatur/Maus
nicht in Ellenbogenebene; (6) keine Steh-Sitz-Dynamik.

\textbf{(b)} Maßnahmen: (1) Bildschirm erhöhen bzw.\ externen Monitor nutzen; (2) Bildschirm
im rechten Winkel zum Fenster, reflexionsarme Oberfläche; (3) höhenverstellbarer Stuhl nach
DIN EN 1335; (4) Sitzhöhe anpassen oder Fußhocker; (5) externe Tastatur/Maus auf
Ellenbogenebene; (6) Steh-Sitz-Tisch und Haltungswechsel.

\textbf{(c)} Ein Laptop vereint Bildschirm und Tastatur. Steht der Bildschirm richtig hoch,
liegt die Tastatur zu hoch --- und umgekehrt. Eine Dockingstation mit externem Bildschirm,
Tastatur und Maus erlaubt die korrekte, getrennte Positionierung jedes Elements.

\textbf{(d)} Dauerhaft gleiche Haltung belastet einseitig Muskulatur und Bandscheiben.
Wechselnde Haltungen (dynamisches Sitzen, Steh-Sitz-Dynamik) entlasten den Rücken,
verbessern die Durchblutung und beugen Verspannungen und Fehlbelastungen vor.
\end{loesung}

\subsection*{Lösung Aufgabe~5}
\begin{loesung}
\textbf{(a)} Dynamisches Sitzen = häufiges Verändern der Sitzposition. Drei Positionen:
\emph{vordere} (Winkel unter 90\textdegree), \emph{mittlere} (aufrecht, 90\textdegree),
\emph{hintere} (zurückgelehnt, größer als 90\textdegree).

\textbf{(b)} 60 \% der Arbeitszeit im Sitzen, 30 \% im Stehen, 10 \% mit gezieltem
Umhergehen.

\textbf{(c)} Zwei bis vier Haltungswechsel pro Stunde; Stehphasen höchstens 20--30 Minuten.

\textbf{(d)} \emph{Synchronmechanik} (Sitz/Lehne fest gekoppelt, Winkel ändert sich beim
Zurücklehnen automatisch), \emph{Wippmechanik} (gekoppelt, Winkel konstant, Schaukel-Gefühl),
\emph{Asynchronmechanik} (Sitz- und Lehnenwinkel unabhängig verstellbar),
\emph{Pendelmechanik} (Steuerung durch Gewichtsverlagerung, wie Sitzball).
\end{loesung}

\subsection*{Lösung Aufgabe~6}
\begin{loesung}
\textbf{(a)} Das \textbf{GS-Zeichen} („Geprüfte Sicherheit") bescheinigt die Einhaltung des
Produktsicherheitsgesetzes und wird \emph{freiwillig} verwendet. Das \textbf{CE-Zeichen} ist
\emph{gesetzlich vorgeschrieben}; der Hersteller bestätigt damit in Eigenverantwortung die
Einhaltung der EU-Anforderungen (technischer Reisepass für den Binnenmarkt). Beide sind die
einzig gesetzlich geregelten Produktsicherheitszeichen in Europa.

\textbf{(b)} Anforderungen (Auswahl): Arbeitsfläche mind.\ 160 $\times$ 80 cm (Call-Center
120 $\times$ 80 cm), Mindesttiefe 80 cm, höheneinstellbar 65--85 cm, abgerundete Kanten,
reflexionsarme Oberfläche, geringe Wärmeleitfähigkeit, genügend Beinfreiheit/Stauraum.
Steh-Sitz-Arbeitstisch: Arbeitshöhe 62--120 cm (Elektromotor oder Gasfeder).

\textbf{(c)} Die ArbStättV legt fest, was der Arbeitgeber beim Einrichten und Betreiben von
Arbeitsstätten für Sicherheit und Gesundheitsschutz beachten muss. Die Anforderungen an
Bildschirmarbeitsplätze stehen im Anhang unter Punkt 6 (ehemals Bildschirmarbeitsverordnung),
konkretisiert durch die ASR.

\textbf{(d)} Vier Anforderungen (Auswahl): ausreichend Raum für wechselnde Arbeitshaltungen;
reflexionsarme/blendfreie Bildschirme und Oberflächen; der Arbeitsaufgabe angepasste
Beleuchtung; getrennte, neigbare Tastatur mit Handballenauflage; flimmer- und verzerrungs\-freies
Bild; individuell einstellbare Helligkeit/Kontrast.
\end{loesung}

\end{document}
